\clearpage
\chapter*{ABSTRAK}
\addcontentsline{toc}{chapter}{ABSTRAK}

\begin{center}
    \begin{singlespace}
        {\large\bfseries \ThesisTitle}

        {\normalfont\normalsize Oleh:}

        {\bfseries \AuthorName}
    \end{singlespace}
\end{center}

\begin{singlespace}
    \small

    Sistem \emph{Retrieval-Augmented Generation} (RAG) harus menemukan informasi yang relevan dari dokumen sebelum model bahasa membentuk jawaban. Tantangan ini semakin besar pada arsitektur \emph{multi-domain} karena peraturan perundang-undangan bidang pendidikan di Indonesia dan \emph{User Manual} produk Netgear memiliki struktur, istilah, dan kebutuhan pencarian yang berbeda. Pemotongan dokumen yang mengabaikan struktur dapat memutus konteks atau menghilangkan identitas sumber, sedangkan strategi pengindeksan dan pembentukan hubungan yang dibuat khusus untuk satu \emph{domain} sulit dipelihara ketika sumber baru ditambahkan. Diperlukan \emph{pipeline} yang menjaga keterlacakan informasi sekaligus dapat diperluas tanpa mengubah kode inti.

    Penelitian ini merancang dan mengimplementasikan \emph{pipeline} konstruksi basis pengetahuan pada sistem Omni-RAG sebagai arsitektur \emph{multi-domain} berbasis \emph{plugin}. Konstruksi basis pengetahuan mencakup tiga proses utama, yaitu membagi dokumen menjadi unit informasi yang lebih kecil dengan mempertahankan konteks, memproyeksikan unit tersebut menjadi indeks pencarian \emph{sparse}, \emph{dense}, atau \emph{hybrid}, serta membentuk entitas dan relasi dalam \emph{knowledge graph} dari \emph{chunk} kanonik yang sama. \emph{Canonical artifact}, \emph{provenance}, \emph{profile}, \emph{manifest}, dan kontrak kompatibilitas digunakan sebagai batas integrasi antarkomponen. Dengan rancangan ini, strategi \emph{chunker}, \emph{indexer}, dan \emph{knowledge graph} dapat dipilih atau ditambahkan melalui \emph{plugin} tanpa mengubah orkestrator inti. Komponen \emph{Context PII}, \emph{guardrail}, dan \emph{history compactor} melengkapi pengelolaan konteks percakapan.

    Evaluasi dilakukan secara \emph{retrieval-only} pada sembilan manual Netgear yang mencakup 1.400 halaman dan 200 pertanyaan dwibahasa, serta pada korpus 255 peraturan perundang-undangan bidang pendidikan di Indonesia dengan 100 pertanyaan. Pada \emph{User Manual}, \emph{structure-aware chunker} mencapai \emph{recall} 0,740 pada \(K=25\), sedangkan \emph{dense indexer} mencapai \emph{recall} 0,830 pada \(K=40\). Pada korpus peraturan, \emph{legal structure-aware chunker} mencapai \emph{recall} 0,64 pada \(K=25\), sedangkan \emph{hybrid indexer} mencapai \emph{recall} 0,67. Hasil tersebut menunjukkan bahwa pemeliharaan struktur dan identitas sumber penting untuk menemukan ketentuan hukum yang relevan, sedangkan strategi pencarian yang berbeda dapat dipilih sesuai karakteristik \emph{domain}. Temuan ini terbatas pada korpus dan evaluasi \emph{retrieval} yang digunakan.

    \textit{\textbf{Kata kunci}:} Omni-RAG, \emph{Retrieval-Augmented Generation}, konstruksi basis pengetahuan, \emph{chunking}, \emph{indexer}, \emph{knowledge graph}, arsitektur \emph{plugin}.
\end{singlespace}

\clearpage
